\chapter{Definición del Problema}
\label{cap:definicion_problema}

Los sistemas de domótica más extendidos hoy en día comparten un mismo modelo de funcionamiento: el hogar se conecta a la nube, y la inteligencia que lo gobierna reside en los servidores del fabricante. Amazon Alexa, Google Home, Tuya Smart, por citar solo los más conocidos, ejecutan el procesamiento de órdenes y la toma de decisiones en centros de datos remotos, dejando al dispositivo local la mera función de terminal de entrada y salida.

Este modelo, que ha democratizado el acceso a la domótica, acarrea problemas estructurales difíciles de ignorar.

\section{Dependencia de Conectividad}

El primero es la dependencia de una conexión permanente a Internet. Si el router se queda sin red, por una avería del operador, un corte programado o simplemente porque el usuario vive en una zona con cobertura deficiente, el hogar inteligente se convierte en un conjunto de objetos inertes. El usuario no puede encender una luz, regular el termostato ni consultar la temperatura de su casa. Esta dependencia es particularmente grave en hogares donde la domótica cubre funciones esenciales, como el control de la calefacción o la ventilación.

Estudios del sector~\cite{cnmc_panel_2025} muestran que una parte relevante de los hogares españoles carece de acceso a Internet de banda ancha, con diferencias significativas entre zonas urbanas y rurales. A esto se suman los cortes eléctricos: según el Consejo de Reguladores Europeos de Energía (CEER), el hogar medio europeo sufre entre 1 y 4 cortes de suministro al año~\cite{ceer_outages_2025}. Durante esos cortes no solo se pierde la conectividad externa, sino que los routers domésticos (que dependen de la corriente) dejan de funcionar por completo.

\section{Exposición de Datos Sensibles}

El segundo problema es la exposición de datos personales. Cada interacción con un asistente de voz, cada cambio de temperatura registrado por un sensor, cada horario de encendido y apagado de luces, viaja a servidores gestionados por terceros. El usuario no tiene control sobre dónde se almacenan esos datos, durante cuánto tiempo se conservan ni con qué fines se utilizan.

Casos sonados no faltan. En 2019, una investigaci'on de Bloomberg revel'o que empleados de Amazon escuchaban grabaciones de Alexa para mejorar el reconocimiento de voz, sin que los usuarios fueran informados~\cite{bloomberg_alexa_2019}. M'as all'a de incidentes concretos, el problema es estructural: el modelo de negocio de los fabricantes se basa en la recolecci'on de datos del usuario, que quedan almacenados en servidores de terceros sin que este tenga control sobre su uso o conservaci'on.

El problema de fondo no es técnico sino de incentivos: mientras la domótica comercial dependa de servidores externos para funcionar, la privacidad del usuario será un elemento negociable en la ecuación económica del fabricante.

\section{Latencia y Capacidad de Respuesta}

El tercer problema es la latencia. Incluso en conexiones de banda ancha, el trayecto de ida y vuelta hasta la nube introduce retardos perceptibles. Una orden tan simple como ``apaga la luz'' puede tardar varios segundos en ejecutarse, lo que resulta frustrante en el uso diario. En aplicaciones que requieren respuesta inmediata, como el control de un ventilador cuando la temperatura sube por encima de un umbral, esos segundos marcan la diferencia.

La revisión sistemática de herramientas y aplicaciones de TinyML realizada por Kallimani et al.~\cite{kallimani_tinyml_2023} documenta cómo el procesamiento local puede reducir la latencia de extremo a extremo en órdenes de magnitud, pasando de segundos a milisegundos para tareas de clasificación y control. Esta reducción no solo beneficia la experiencia del usuario, sino que también tiene un impacto energético positivo: un sistema que responde en milisegundos puede reaccionar antes a cambios ambientales, reduciendo el tiempo de funcionamiento innecesario de los actuadores.

\section{Heterogeneidad de Protocolos}

A los problemas anteriores se suma la fragmentación del ecosistema IoT. Los estándares de comunicación son múltiples y raramente interoperan de forma nativa: Tuya, Zigbee, Z-Wave, Bluetooth Low Energy, bus LIN para climatización~\cite{iso_17987_2025}. Cada fabricante apuesta por su propio protocolo, y el usuario que quiere integrar dispositivos de distintas marcas se encuentra con que no hablan entre sí. Las plataformas comerciales resuelven esto llevándolo todo a la nube, pero eso reintroduce los problemas de dependencia y privacidad.

Esta heterogeneidad no es un accidente: es una estrategia comercial. Los ecosistemas cerrados crean barreras de salida para el usuario, que no puede mezclar dispositivos de distintos fabricantes sin perder funcionalidades~\cite{home_assistant_2019}.

\section{Restricción Económica como Factor de Diseño}

Por último, el proyecto opera bajo una restricción económica real. El autor es estudiante sin ingresos, y cada dispositivo ha sido seleccionado priorizando el coste mínimo frente a la comodidad de integración. Un analizador de protocolos BLE (como el nRF52840, con un coste superior a 50~\euro{}) o un mando universal por infrarrojos habrían acelerado el desarrollo, pero se ha optado por la ingeniería inversa con herramientas ya disponibles.

Esta restricción no es una limitación del proyecto, sino una decisión deliberada: demostrar que la domótica inteligente puede construirse con dispositivos de consumo masivo y presupuestos ajustados. El coste total en licencias de software y servicios cloud ha sido de 0~\euro{}, gracias a que todos los dispositivos o bien ya estaban en posesión del autor, o bien fueron adquiridos con anterioridad por motivos ajenos al proyecto.

\section{Síntesis del Problema}

Todos estos problemas (dependencia de red, exposición de datos, latencia, fragmentación de protocolos y restricción económica) comparten una misma raíz: la delegación de la inteligencia del hogar a infraestructuras externas. Este proyecto propone una respuesta directa: un ecosistema de domótica diseñado desde el primer día para operar \textbf{completamente} en la red local, sin depender de la nube para ninguna función crítica, y capaz de integrar hardware heterogéneo mediante una arquitectura de software desacoplada y orquestada por inteligencia artificial local.
